\chapter{Smooth support and multiplicity}\label{sec:varifold-constancy}
Stationarity constrains both geometry and multiplicity. On a connected smooth support it forces constant multiplicity and minimality. Without stationarity a density can vary, and its derivative contributes to first variation. The main source is \cite{Minter}, Lemma 0.7 and its proof, pp.~8--9.

\begin{lemma}[A distributionally constant function is constant]\label{lem:weak-gradient-zero}
Let $W\subset\mathbb R^m$ be connected and open, and let $h\in L^1_{\mathrm{loc}}(W)$. If
\[
 \int_W h\,\partial_i\phi\,dx=0
 \quad(\phi\in C_c^\infty(W),\ i=1,\ldots,m),
\]
then $h$ is equal almost everywhere to one constant.
\end{lemma}
\begin{proof}
On a ball compactly contained in $W$, mollification gives smooth functions whose derivatives vanish on smaller concentric balls. Each mollification is constant on each such smaller ball. Its $L^1$ convergence shows that the limiting function is almost everywhere constant there as well: the constants converge, as seen by taking their averages. Overlapping balls have the same constant. Any two points of a connected open subset of Euclidean space can be joined by a finite chain of overlapping balls, so the local constants agree. This supplies the distributional step in \cite{Minter}, Lemma 0.7.
\end{proof}

\begin{proposition}[Tangential stationarity forces constant density]\label{prop:tangential-density-constancy}
Let $M$ be a connected $C^2$ Riemannian manifold without boundary and let $\theta\in L^1_{\mathrm{loc}}(M)$. If
\[
 \int_M\theta\operatorname{div}_M Y\,dA=0
 \quad\text{for all compactly supported }C^1\text{ tangent fields }Y,
\]
then $\theta$ is constant almost everywhere.
\end{proposition}
\begin{proof}
In coordinates with volume element $\sqrt{g}\,dx$, the divergence is
$\operatorname{div}_M Y=(\sqrt g)^{-1}\partial_i(\sqrt gY^i)$. Choose $Y^i=\phi/\sqrt g$ for one component and zero for the others. The hypothesis becomes $\int\theta\partial_i\phi\,dx=0$. Lemma \ref{lem:weak-gradient-zero} gives a constant on each connected chart. Constants agree on overlaps and hence throughout connected $M$. The volume factor must be included in this coordinate argument; it is absorbed into the chosen test field. Compare \cite{Minter}, the second stage of the proof of Lemma 0.7.
\end{proof}

\begin{theorem}[Constancy on a smooth connected support]\label{thm:varifold-smooth-constancy}
Let $V$ be a stationary integral $m$-varifold in $U\subset\mathbb R^N$. Suppose its support is contained in a connected, relatively closed, embedded $C^2$ $m$-submanifold $M\subset U$ without boundary in $U$. Then
\[
 V=q|M|\quad\text{in }U
\]
for an integer $q\geq0$. If $V\ne0$, then $q\geq1$, $M$ is minimal, and $\operatorname{spt}\|V\|=M$ in $U$. This is the local smooth-support constancy theorem of \cite{Minter}, Lemma 0.7, with the connected portion and ambient open set specified.
\end{theorem}
\begin{proof}
Integrality gives a representation on a rectifiable set with integer density. Since the weight is supported on $M$, it is absolutely continuous with respect to $\mathcal H^m\llcorner M$; extend the density by zero on the unused part of $M$. Call it $\theta\geq0$.

We also identify the plane component. Locally choose $C^2$ defining functions $f_\alpha$ for $M$ whose gradients span its normal space, and a nonnegative cutoff $\eta$ supported in that coordinate neighborhood. The compactly supported field $X=\eta f_\alpha\nabla f_\alpha$ satisfies on $M$
\[
 \operatorname{div}_S X=\eta|\operatorname{proj}_S\nabla f_\alpha|^2.
\]
Stationarity forces this nonnegative integral to vanish. A countable cover and cutoffs positive on smaller neighborhoods show that almost every varifold plane is orthogonal to every normal gradient. Equality of dimensions implies $S=T_xM$. Thus $V=V(M,\theta)$.

Every compactly supported $C^1$ tangent field on $M$ extends locally to an ambient test field; a partition of unity gives a compactly supported extension in $U$. Stationarity and Proposition \ref{prop:tangential-density-constancy} imply $\theta=q$ almost everywhere on connected $M$. Integrality makes $q$ a nonnegative integer. If it is positive, Proposition \ref{prop:weighted-smooth-first-variation} gives
$0=\delta V(X)=-q\int_M\mathbf H_M\cdot X\,dA$. Hence the continuous mean-curvature vector vanishes. A positive constant density gives positive mass in every neighborhood of every point of $M$, so its relative support is all of $M$. The defining-function argument follows \cite{Minter}, pp.~8--9.
\end{proof}

\begin{example}[A multiplicity jump produces first variation]\label{ex:jump-multiplicity}
On $P=\mathbb R^m\times\{0\}\subset\mathbb R^{m+1}$, set $\theta=1$ when $x_1<0$ and $\theta=2$ when $x_1>0$. This defines an integral varifold with smooth planar support, but for a compactly supported ambient field $X$,
\[
 \delta V(X)=-\int_{P\cap\{x_1=0\}}X\cdot e_1\,d\mathcal H^{m-1}.
\]
Indeed the two half-planes have opposite outward conormals, and their weights leave coefficient $1-2=-1$ on the interface. Thus the varifold is not stationary. Smoothness of support alone does not imply constant multiplicity, and the singular first variation is not represented by a function integrable with respect to the $m$-dimensional weight.
\end{example}

\begin{proposition}[First variation of a smooth positive density]\label{prop:weighted-smooth-first-variation}
Let $M\subset U$ be a relatively closed embedded $C^2$ submanifold without boundary in $U$, and let $\theta>0$ be $C^1$ on $M$. Assume the associated weight is locally finite. Then
\[
 \delta V(M,\theta)(X)
 =-\int_M(\mathbf H_M+\nabla^M\log\theta)\cdot X\,\theta\,dA.
\]
Here $\mathbf H_M$ is the geometric mean-curvature vector, equal to $\Delta_M F$ for the inclusion map. The tangential term records variation of the weight.
\end{proposition}
\begin{proof}
Decompose $X=X^T+X^\perp$. The smooth divergence identity is
$\operatorname{div}_M X=\operatorname{div}_M X^T-\mathbf H_M\cdot X$. Integration by parts against $\theta$ converts the first term to $-\nabla^M\theta\cdot X$. Factor out $\theta$ to obtain the formula. Relative closedness ensures that a compact ambient support meets $M$ in a compact set, so no boundary term occurs. This computes Definition \ref{def:first-variation} explicitly and complements the stationarity calculation in \cite{Minter}, Lemma 0.7.
\end{proof}

\begin{example}[Generalized mean curvature can have a tangential part]\label{ex:weighted-plane-mean-curvature}
On the same plane $P$, use the smooth positive density $\theta(x)=e^{x_1}$. The weight is locally finite and the generalized mean curvature in Definition \ref{def:first-variation} is $e_1$, entirely tangent to $P$. This varifold is rectifiable but not integral. Thus normality of generalized mean curvature is not a consequence of that definition for every rectifiable varifold; additional hypotheses, such as those in the integral-varifold theory, matter.
\end{example}
